\booksection{Anexo del Capítulo 11: Tecnologia de imagen profesional}

\chapterintro{Para comprender verdaderamente los secretos de la fotografía y video móvil, necesitas conocimientos técnicos de imagen profesionales. Esto no solo te ayudará a usar mejor la cámara del móvil, sino también a tomar mejores decisiones al elegir equipo.}

\booksubsection{Resolución, tasa de frames y bitrate: los tres elementos de la calidad de video}

\subsubsection*{Resolución:}

La resolución determina la claridad de la imagen; se expresa en píxeles.

\begin{itemize}
  \item \textbf{720P (HD)}: 1280×720 píxeles; calidad de nivel principiante.
  \item \textbf{1080P (Full HD)}: 1920×1080 píxeles; la resolución más popular actualmente.
  \item \textbf{2K (QHD)}: 2560×1440 píxeles; más nítida que 1080P.
  \item \textbf{4K (UHD)}: 3840×2160 píxeles; ultra alta definición, rico en detalles.
  \item \textbf{8K}: 7680×4320 píxeles; estándar del futuro, soporte limitado en dispositivos actuales.
\end{itemize}

Mayor resolución = imagen más clara, pero archivos más grandes y mayor demanda de almacenamiento y potencia de procesamiento.

\subsubsection*{Tasa de frames (Frame Rate):}

La tasa de frames determina la fluidez; se mide en frames por segúndo (fps).

\begin{itemize}
  \item \textbf{24fps}: Tasa de frames estándar de cine; tiene una única "sensación cinematográfica".
  \item \textbf{30fps}: Común en TV y video online; fluido y natural.
  \item \textbf{60fps}: Alta tasa de frames; ideal para escenas de acción y grabación de juegos; muy fluido.
  \item \textbf{120fps/240fps}: Súper alta tasa de frames; usable para video en cámara lenta.
\end{itemize}

Más fps = más fluidez, mejor efecto de ralentización, pero archivos más grandes.

\subsubsection*{Bitrate:}

El bitrate determina la densidad de datos del video; se expresa en cantidad de datos por segúndo (Mbps).

\begin{itemize}
  \item Mayor bitrate = mejor calidad, más detalles, pero archivos más grandes.
  \item Bitrate muy bajo produce artefactos de compresión como mosaicos y bloques de color.
  \item La grabación profesional suele requerir alto bitrate para preservar margen para postproducción.
\end{itemize}

\booksubsection{Formatos de codificación de video: H.264, H.265 y ProRes}

\subsubsection*{H.264 (AVC):}

\begin{itemize}
  \item \textbf{Características}:El formato de codificación más usado; compatibilidad excelente.
  \item \textbf{Ventajas}:Casi todos los dispositivos lo reproducen; eficiencia de compresión moderada.
  \item \textbf{Desventajas}:Comparado con formatos nuevos, menor eficiencia de compresión; archivos más grandes.
  \item \textbf{Casos de uso}:Compartir video en red; grabación diaria de usuarios normales.
\end{itemize}

\subsubsection*{H.265 (HEVC):}

\begin{itemize}
  \item \textbf{Características}:Versión mejorada de H.264; eficiencia de compresión un \textasciitilde{}50\% superior.
  \item \textbf{Ventajas}:Misma calidad con archivos menores; soporta mayor resolución y tasa de frames.
  \item \textbf{Desventajas}:Codificación/decodificación requiere procesador más potente; algúnos dispositivos antiguos incompatibles.
  \item \textbf{Casos de uso}:Grabación de video 4K; situaciones con almacenamiento limitado.
\end{itemize}

\subsubsection*{Apple ProRes:}

\begin{itemize}
  \item \textbf{Características}: Códec de video profesional desarrollado por Apple; compresión casi sin pérdida.
  \item \textbf{Ventajas}: Calidad de imagen extremadamente alta; amplio margen para postproducción; color preciso.
  \item \textbf{Desventajas}: Archivos enormes; requiere mucho espacio de almacenamiento; necesita equipo profesional para procesar.
  \item \textbf{Casos de uso}: Producción de video profesional; grabación de nivel cinematográfico.
\end{itemize}

\subsubsection*{ProRes RAW:}

\begin{itemize}
  \item \textbf{Características}: Combina las ventajas de los formatos ProRes y RAW.
  \item \textbf{¿Qué es RAW?}: Datos brutos del sensor sin procesar; conserva la máxima cantidad de información.
  \item \textbf{Ventajas}: Amplísimo margen de ajuste en postproducción; puedes cambiar exposición, balance de blancos y otros parámetros sin perder calidad.
  \item \textbf{Desventajas}: Archivos inmensamente grandes; requieren software de postproducción profesional.
  \item \textbf{Casos de uso}: Produccion profesional de alta gama. Desde iPhone 15 Pro se soporta grabacion de video Apple ProRes (hasta 4K 60fps); iPhone 17 Pro anade ProRes RAW grabado directamente en el dispositivo.
\end{itemize}

\booksubsection{HDR: tecnología de alto rango dinámico}

La tecnología HDR (High Dynamic Range) puede conservar simultáneamente los detalles de las partes más claras y más oscuras de la imagen.

\begin{itemize}
  \item \textbf{HDR10}: El estándar HDR más básico; buena compatibilidad.
  \item \textbf{HDR10+}: Estándar promovido por Samsung; soporta metadatos dinámicos.
  \item \textbf{Dolby Vision}: El estándar HDR de máxima gama desarrollado por Dolby; la mejor calidad de imagen, pero requiere licencia.
  \item \textbf{HLG}: Estándar HDR de nivel broadcast; compatible con monitores normales.
\end{itemize}

Al grabar video HDR, necesitas un monitor compatible con HDR para ver el efecto completo.

\booksubsection{Espacio de color y profundidad de color}

\subsubsection*{Espacio de color:}

\begin{itemize}
  \item \textbf{sRGB}: Espacio de color estándar; la mejor compatibilidad.
  \item \textbf{Display P3}: Estándar de gamut amplio promovido por Apple; colores más ricos.
  \item \textbf{Rec.2020}: Espacio de color estándar para video HDR; el gamut más amplio.
\end{itemize}

\subsubsection*{Profundidad de color:}

\begin{itemize}
  \item \textbf{8-bit}: 16,7 millones de colores; estándar de video normal.
  \item \textbf{10-bit}: 1.070 millones de colores; estándar de video profesional; transiciones de color más suaves.
  \item \textbf{12-bit}: 68.000 millones de colores; estándar profesional de alta gama.
\end{itemize}

\begin{tcbinfo}{Consejos prácticos}
\begin{itemize}
  \item Grabación diaria: 1080P/4K + 30fps + codificación H.265; equilibra calidad y tamaño de archivo.
  \item Grabación de acción: 1080P/4K + 60fps; garantiza fluidez.
  \item Creación profesional: 4K + ProRes (si es compatible); conserva el máximo margen para postproducción.
  \item Almacenamiento: los videos de alto bitrate ocupan mucho espacio; se recomienda hacer copias de seguridad periódicas al ordenador o a la nube.
\end{itemize}

\end{tcbinfo}

\begin{tcbpractice}{Explora: ajustes de video de tu móvil}
\begin{enumerate}
  \item Abre la App de cámara, cambia al modo video.
  \item Revisa las opciones de resolución y tasa de frames; prueba diferentes combinaciones.
  \item Si tu móvil lo soporta, verifica si hay opciones avanzadas como HDR o ProRes.
  \item Graba varios clips con distintos ajustes; compara calidad de imagen y tamaño de archivo.
\end{enumerate}

\end{tcbpractice}

\booksubsection{Comparacion de formato RAW, HEIF y JPEG}

\subsubsection*{Que es el formato RAW:}

RAW son los “datos en bruto” que registra el sensor de la camara: informacion pura sin compresion ni procesado. Es como un ingrediente crudo que hay que “cocinar” en posproduccion. Como el telefono tiene almacenamiento y capacidad de calculo limitados, por defecto genera JPEG o HEIF ya “cocinados”; RAW hay que activarlo expresamente.

\subsubsection*{Comparacion de RAW, HEIF y JPEG:}

\begin{itemize}
  \item \textbf{JPEG}: El formato mas antiguo y con mejor compatibilidad (cualquier dispositivo lo abre). 8-bit de profundidad de color, compresion con perdida, archivo mas pequeno (2-5 MB por foto). Espacio de posproduccion limitado; forzar las sombras produce artefactos de color.
  \item \textbf{HEIF/HEIC}: Nuevo formato impulsado por Apple; con la misma calidad, el archivo es la mitad que JPEG (1-3 MB por foto). Soporta 10-bit de color y canal alfa. El inconveniente es que algunos dispositivos antiguos y Windows necesitan un plugin para abrirlo.
  \item \textbf{RAW}: Datos en bruto sin compresion, el archivo mas grande (20-100 MB por foto). Profundidad de color de 12-16 bit, maximo espacio de posproduccion. El inconveniente es que requiere software profesional (Lightroom, Snapseed, etc.) para abrir y procesar, no se puede compartir directamente.
\end{itemize}

\subsubsection*{Por que los fotografos profesionales disparan en RAW:}

\begin{itemize}
  \item \textbf{Ajuste del balance de blancos en posproduccion sin perdida}: Si el balance de blancos del JPEG es erroneo (por ejemplo, amarillento bajo luz tungsteno), corregirlo en posproduccion reduce la calidad; en RAW se puede reajustar la temperatura de color sin perdida alguna.
  \item \textbf{Mayor tolerancia de exposicion}: En JPEG, las zonas subexpuestas 2 pasos son negro absoluto; al aclararlas aparecen ruido. En RAW se pueden recuperar detalles de zonas subexpuestas 3-4 pasos. Lo mismo aplica a las sobreexpuestas.
  \item \textbf{Mas detalle en sombras y altas luces}: El rango dinamico de RAW es muy superior al de JPEG; hay mucho mas margen para que las altas luces no se quemen y las sombras no sean negro absoluto.
  \item \textbf{Nitidez, reduccion de ruido y correccion de lente controlables en posproduccion}: El procesado interno de la camara es "a fondo perdido"; en RAW se puede ajustar repetidamente partiendo siempre de los datos originales.
\end{itemize}

\subsubsection*{Como activar RAW en el telefono:}

\begin{itemize}
  \item \textbf{iPhone ProRAW}: Solo lo soportan iPhone 12 Pro y modelos Pro superiores. Abre "Ajustes > Camara > Formato" y activa "Apple ProRAW". Aparecera un boton "RAW" en la parte superior de la interfaz de la camara; tocalo para activarlo. Cada foto ocupa unos 25 MB; activa solo cuando lo necesites.
  \item \textbf{Android Motion Raw / Ultra RAW}: Cada marca lo llama distinto. Xiaomi/OPPO/vivo en gamas altas lo encuentran en "Modo Pro" o "Mas"; Huawei lo activa en "Modo Pro"; Samsung Galaxy S lo ofrece en modo "Pro". El tamano del archivo suele ser de 20-50 MB.
  \item \textbf{Recomendacion general}: Para fotos cotidianas no hace falta activar RAW; el archivo es demasiado grande y ocupa espacio. Activar solo en escenas importantes (paisaje, retrato, cuando se necesita posproduccion).
\end{itemize}

\booksubsection{Principios de la fotografia computacional}

La fotografia movil moderna no puede prescindir de la “fotografia computacional”, es decir, generar imagenes mediante algoritmos y no solo con lentes y sensores. Comprender estos principios basicos te ayuda a saber que hace el telefono en el instante de pulsar el disparador.

\subsubsection*{Fusion de multiples fotogramas:}

Al pulsar el disparador, el telefono en realidad capta varias fotos. El algoritmo las alinea y las fusiona ponderadamente para obtener una imagen final de alta calidad.

\begin{itemize}
  \item \textbf{Modo nocturno}: En 1-3 segundos dispara varias fotos subexpuestas, las alinea y las superpone: el ruido se cancela mutuamente y la luminosidad se acumula. Asi se consiguen imagenes limpias incluso con poca luz.
  \item \textbf{HDR (alto rango dinamico)}: Captura tres (o mas) fotos subexpuesta, normal y sobreexpuesta, y las fusiona en una sola imagen con detalle en altas luces y sombras. Muy eficaz en escenas a contraluz.
  \item \textbf{Deep Fusion (Apple) / Photonic Engine (motor de fotones)}: Para escenas con luz media o baja, analiza multiples fotogramas pixel a pixel y los fusiona; mejora sobre todo texturas y detalles: sueteres, pelo y similares se ven claramente mas nítidos.
\end{itemize}

\subsubsection*{Segmentacion semantica:}

La IA no procesa toda la imagen por igual; primero “comprende” que hay en la escena (cielo, rostros, cesped, edificios, ropa) y luego aplica distintas estrategias a cada zona.

\begin{itemize}
  \item \textbf{Cielo}: Refuerza el azul por separado y reduce el brillo para evitar sobreexposicion.
  \item \textbf{Rostros}: Detecta zonas de tono piel, las aclara y suaviza sin tocar los rasos faciales ni el detalle del pelo.
  \item \textbf{Cesped / hojas}: Aumenta la saturacion del verde y mejora la nitidez del detalle.
  \item \textbf{Desenfoque del fondo}: En modo retrato, la IA detecta el contorno del sujeto y desenfoca el fondo manteniendo el sujeto nítido.
\end{itemize}

\subsubsection*{Aplicaciones del aprendizaje profundo:}

\begin{itemize}
  \item \textbf{Reconocimiento de escena por IA}: La camara reconoce automaticamente escenas como "comida", "mascota", "noche", "documento" y optimiza los parametros correspondientes. Por ejemplo, el modo comida intensifica los tonos calidos y la saturacion para que los alimentos parezcan mas apetitosos.
  \item \textbf{Desenfoque de retrato por IA}: Las lentes del telefono no pueden lograr fisicamente la poca profundidad de campo de una camara reflex; se usa aprendizaje profundo para predecir los bordes del sujeto y luego desenfocar el fondo. Los algoritmos mas recientes ya manejan bien detalles como mechones de pelo o bordes de cristal.
  \item \textbf{Eliminacion por IA / Expansion de imagen por IA}: Selecciona un transeunte o un objeto molesto y la IA "rellena" la zona tapada a partir del contexto. Es una aplicacion de IA generativa en el ambito de la imagen.
\end{itemize}

\subsubsection*{Fotografia computacional frente a fotografia optica:}

\begin{itemize}
  \item \textbf{Ventajas de la fotografia computacional}: Con un tamano reducido se obtiene buena calidad; tambien funciona en condiciones extremas como poca luz o contraluz; el usuario medio puede sacar fotos "bonitas" sin necesidad de conocimientos fotograficos.
  \item \textbf{Inconvenientes de la fotografia computacional}: A veces la imagen esta "sobreprocesada": nitidez excesiva, color distorsionado, cielos poco naturales; la fusion de multiples fotogramas deja fantasmas si el sujeto se mueve; los detalles estan "calculados", no registrados de forma real.
  \item \textbf{Ventajas de la fotografia optica}: Sensores grandes con lentes de gran apertura captan la luz de forma real, calidad natural y amplio margen de posproduccion; no genera artefactos por "alucinacion" del algoritmo.
  \item \textbf{Tendencia}: Los telefonos de gama alta actuales combinan "optica + computacion": sensor grande para captar informacion real y algoritmos para optimizar el resultado. No son opuestos.
\end{itemize}

\booksubsection{Tipos de lente y distancias focales}

Los telefonos de gama alta suelen tener tres o cuatro lentes, cada uno con una distancia focal distinta. Entender las distancias focales te ayuda a elegir el lente adecuado para cada escena.

\subsubsection*{Lentes comunes y distancia focal equivalente:}

\begin{itemize}
  \item \textbf{Ultra gran angular (equiv. 13-24 mm)}: Campo de vision mas amplio; ideal para paisajes, arquitectura y espacios interiores. Capta grandes escenas, pero los bordes sufren distorsion (las lineas rectas se curvan) y los rostros se estiran en retratos. Adecuado para "escenas amplias", no para "primeros planos".
  \item \textbf{Camara principal (equiv. 24-28 mm)}: El lente mas usado a diario, con la mejor calidad y la mayor apertura. Equilibra campo de vision y calidad; sirve para personas y paisajes. El "1x" de la mayoria de telefonos es la camara principal.
  \item \textbf{Teleobjetivo (equiv. 70-200 mm)}: Para paisajes lejanos y primeros planos de retrato. El teleobjetivo comprime la sensacion de profundidad (las montanas distantes y la persona cercana parecen mas juntas) y en retratos ofrece proporciones faciales mas naturales y mayor desenfoque del fondo.
  \item \textbf{Teleobjetivo periscopico (equiv. 200 mm+)}: Mediante un diseno de ruta optica reflejada, logra una distancia focal mas larga dentro del grosor del telefono. Capta detalles a gran distancia; ideal para conciertos, eventos deportivos y fauna. Solo en gamas altas.
\end{itemize}

\subsubsection*{Compromiso entre tamano de sensor y tamano de pixel:}

El sensor es la “base” y determina cuanta luz se capta. En el sector se dice que “un sensor mas grande lo aplasta todo”: con el mismo numero de pixeles, cuanto mayor sea el sensor, mas grandes seran los pixeles individuales, mas luz entrara y mejor sera la calidad de imagen.

\begin{itemize}
  \item \textbf{Sensor de gran formato}: Captura mucha luz, buena calidad nocturna y amplio rango dinamico. El inconveniente es que el telefono debe ser mas grueso (el modulo de camara sobresale) y el coste es mayor.
  \item \textbf{Sensor pequeno de alta resolucion}: Con buena luz diurna, los muchos megapixeles producen fotos mas nítidas y facilitan el recorte posterior. Pero con poca luz, cada pixel recibe menos luz y aparece ruido.
  \item \textbf{Eleccion en gamas altas}: La camara principal suele ser de 1/1,3 pulgadas a 1 pulgada (algunos modelos ya usan 1 pulgada); la resolucion ronda los 50 megapixeles: el punto de equilibrio entre calidad y tamano.
\end{itemize}

\subsubsection*{Tecnologia de fusion de pixeles (Pixel Binning):}

Para combinar alta resolucion de dia con pixeles grandes de noche, los fabricantes usan “fusion de pixeles”: fusiona varios pixeles pequenos en uno mas grande para la salida.

\begin{itemize}
  \item \textbf{Fusion cuadruple (Quad Bayer)}: Sensor de 50 MP; genera fotos de 12,5 MP. Cada pixel de salida se compone a partir de 4 pixeles adyacentes; la luz captada aumenta 4 veces y la imagen es mas limpia.
  \item \textbf{Fusion nonagonal (Nonacell)}: Para sensores de mas de 100 MP; genera aprox. 11-12 MP. Mismo principio, con mayor entrada de luz.
  \item \textbf{Efecto real}: En el modo de foto por defecto, obtienes la imagen fusionada de menor resolucion: buena calidad y archivo pequeno. Cuando necesitas maximo detalle (documentos, paisajes lejanos), cambia manualmente a "modo de alta resolucion" para obtener la resolucion completa.
\end{itemize}

\booksubsection{Modos de captura especiales}

Ademas de las fotos cotidianas, el telefono tiene varios modos especiales que, usados bien, permiten captar imagenes que de otro modo serian imposibles.

\subsubsection*{Modo panoramico:}

\begin{itemize}
  \item Al disparar, mueve el telefono lentamente en horizontal siguiendo la flecha; no vayas rapido o habra errores de alineacion.
  \item Mantén el telefono estable en vertical y desplazalo solo en horizontal. Si oscila arriba y abajo, el horizonte ondulara.
  \item Usalo con precaucion si hay personas o coches en movimiento: aparecera "medio cuerpo" o "fantasmas".
  \item El modo panoramico tambien se puede usar en vertical: disparo vertical y panoramica horizontal; ideal para edificios altos o cascadas.
\end{itemize}

\subsubsection*{Camara lenta y lapso de tiempo:}

\begin{itemize}
  \item \textbf{Camara lenta}: Graba a 120 fps o 240 fps y la reproduccion se vuelve camara lenta. Ideal para gotas de agua, momentos deportivos o saltos de mascotas. Necesita buena luz; en modo lento el tiempo de exposicion es corto y con poca luz la calidad cae.
  \item \textbf{Lapso de tiempo}: Toma varias fotos a intervalos y genera un video acelerado. Ideal para trafico, movimiento de nubes, atardeceres o apertura de flores. La duracion va de minutos a horas; usa un tripode para fijar el telefono, si lo sujetas con la mano saldra movido.
  \item \textbf{Truco}: Al grabar nubes en lapso de tiempo, coloca el telefono en un soporte y graba 15-30 minutos; obtendras unos 10 segundos de video con resultados impresionantes.
\end{itemize}

\subsubsection*{Modo cielo estrellado / astronomico:}

\begin{itemize}
  \item Hay que alejarse de la contaminacion luminica de la ciudad y buscar las afueras o la montana sin farolas. En la ciudad no se capta la Via Lactea.
  \item Imprescindible tripode; con la mano saldra movido. El telefono debe estar firme de varios segundos a decenas de segundos.
  \item Cielo despejado y sin nubes; la luna no debe ser demasiado grande (luna llena oculta las estrellas). Las noches oscuras de principios o finales de mes lunar son las mas adecuadas.
  \item El "Modo nocturno" del iPhone reconoce el cielo estrellado y alarga la exposicion (hasta 30 s); los modelos Android de gama alta tienen "Modo estrella" o "Modo astronomico" con resultados similares. Ten paciencia: tras pulsar el disparador, no muevas el telefono.
\end{itemize}

\subsubsection*{Modo macro:}

\begin{itemize}
  \item Para fotografiar objetos pequenos como flores, insectos, joyas o primeros planos de texto. Puede enfocar a 1-2 cm de distancia.
  \item : Some phones have a dedicated macro lens (usually a 2 MP or 5 MP small-sensor lens) with average quality; flagship phones use the ultra-wide lens to double as macro, with better quality.
  \item : When shooting close-up, the lens is very close to the subject and may block light; use a desk lamp or flashlight to supplement lighting if needed.
  \item : The focus distance is extremely short and the depth of field is very shallow; moving the subject a few millimeters forward or back may cause loss of focus. Keep your hand steady and ideally keep the subject still.
\end{itemize}

\begin{tcbinfo}{Practical tips for still photography}
\begin{itemize}
  \item : For everyday shooting: the default HEIF/JPEG is fine; small file and sufficient quality.
  \item : For important scenes (landscape, portrait, post-processing needed): enable RAW to preserve maximum post-processing flexibility.
  \item : Choose the right lens: ultra-wide for landscape, telephoto for portraits, main camera for everyday use.
  \item : For night scenes and starry skies: always secure the phone; enable night mode or astrophotography mode.
  \item : Large-sensor phones (e.g., 1-inch sensor models) have a higher quality ceiling, but good operating habits remain key.
\end{itemize}

\end{tcbinfo}

\begin{tcbpractice}{Hands-on: get to know your camera capabilities}
\begin{enumerate}
  \item Comprueba si tu telefono soporta disparo en RAW (Ajustes > Camara > Formato), activa la opcion y dispara una foto; compara el tamano del archivo JPEG y del RAW en la misma escena.
  \item Cambia entre ultra gran angular, camara principal y teleobjetivo; dispara el mismo paisaje lejano con cada uno para notar la diferencia de angulo de vision y de calidad.
  \item Busca una escena con modo nocturno o modo estrella, fija el telefono y dispara una foto nocturna.
  \item Prueba a disparar un panorama horizontal de una escena amplia; mantén el telefono estable mientras lo mueves.
\end{enumerate}

\end{tcbpractice}
